% Methods (synced to P0 round-17 state, 2026-09-09)
\section{Methods}
\label{sec:methods}

\subsection{Cohort and endpoint definitions}
\label{ssec:methods_cohort}

We analyzed human platelet single-cell transcriptomes from public cohorts
(Table~\ref{tab:cohort}): 240 COVID-19 patients (25 fatal), 103 sepsis
patients in the dataset (19 fatal), and 70 healthy donors retained as a
reference group only. Multiple draws of the same patient (longitudinal samples
or replicate libraries) were merged into a single patient-level unit
(\texttt{patient\_id}) before any analysis, so that all of a patient's cells
and draws share one cross-validation group and one aggregated score.
GSE151263 (7 patients) does not report per-patient outcomes in the source
record and is excluded from all mortality analyses; it is retained in the
dataset for completeness. This eligibility decision follows from source-record
availability and was fixed at dataset assembly, independently of any
outcome-associated result. No reported table or figure contains GSE151263
patients; the cohort appears only in the dataset description
(Table~\ref{tab:cohort}). In arithmetic terms: of the 103 sepsis patients in
the dataset, 7 enter no outcome-labeled analysis, 91 (GSE279452) form the
main sepsis arm, and all 96 outcome-labeled patients (91 plus 5 from
GSE167363) form the reverse-transfer training pool.
The primary endpoint was mortality: in-hospital
death for the COVID-19 cohorts and GSE167363, and 28-day death for GSE279452;
the difference is stated wherever the GSE279452 arm is used. Survivors
include patients discharged alive at any severity grade (severe cases are
ICU-discharged survivors, labeled as the negative class). Outcome labels were
verified against the source records for every fatal patient
(Supplementary Note~\ref{sn:label_provenance}).

\subsection{Foundation models and unified input}
\label{ssec:methods_models}

We evaluated five single-cell foundation models (scFMs) chosen to span the
main representation families (Table~\ref{tab:models}): \model{STATE},
\model{UCE}, \model{Geneformer}, \model{scPRINT-2}, and \model{scGPT}. To
remove input preprocessing as a confounder, every model received the identical
raw AnnData object and applied its own native preprocessing pipeline as
recommended by its authors; embeddings were then used frozen, with no
fine-tuning, generative, or task-specific capabilities of any model. For
\model{STATE}, a two-component system, we evaluate the state-embedding
component (SE-600M): the state-transition module is trained for
perturbation-response prediction and has no natural mapping to a
patient-level discrimination task, and SE is the system's representation
component. Platelets are anucleate cell fragments carrying
megakaryocyte-derived mRNA; in the three pretraining corpora whose
composition can be numerically bounded they account for approximately
0.12\% of cells, and \model{STATE}'s corpus is bounded lower still (below
0.05\%) through its public component documentation, while only
\model{Geneformer} publishes no cell-type-level corpus composition at all
(Table~\ref{tab:models} and Supplementary Note~\ref{sn:corpus_platelet}).
They are therefore an underrepresented substrate in every corpus that can
be checked --- documented as absent from some, bounded near one per mille
in the rest.

\subsection{Evaluation protocol}
\label{ssec:methods_cv}

The primary evaluation unit is the patient. Classifiers were trained and
evaluated under repeated patient-level cross-validation: 10 repeats of 5-fold
\texttt{StratifiedGroupKFold} (\texttt{shuffle=True}, seeds 42--51) with
\texttt{patient\_id} as the group variable, so that no patient's cells or
draws ever span training and test folds. Within each repeat, out-of-fold cell
predictions were averaged per patient (mean; median and trimmed-mean as
sensitivity analyses, Supplementary Note~\ref{sn:aggregation}) to yield one
score per patient. The primary metric is patient-level out-of-fold ROC-AUC
(patient-level AUC) computed as the repeat-mean, reported with the standard
deviation across repeats (repeat-sd) as a measure of split-to-split stability.

We pre-specified a single primary inferential question: does the best
foundation model beat the from-scratch baseline? Because quoting the best of
five correlated candidates inflates the error rate, this contrast was tested
with a winner-corrected max-statistic permutation test: mortality labels were
permuted at the patient level (10{,}000 permutations) on the frozen
seed-ensemble scores, and the familywise $P$ is the fraction of permutations
whose maximum foundation-model-minus-baseline AUC difference reaches the
observed maximum; a recentered bootstrap is reported as a cross-check
(Supplementary Note~\ref{sn:winner_test}).
Every other contrast is reported as
exploratory, with nominal $p$-values from exact DeLong's test labeled as such
(paired-bootstrap agreement documented in Supplementary
Note~\ref{sn:se_check}), and the stability of the primary result was assessed by
leave-one-fatal-patient-out refits (Supplementary Note~\ref{sn:fragility}).
Uncertainty quantification was matched to each analysis
design: repeated-CV for fold-based within-disease evaluation, patient-clustered
bootstrap (200 iterations) for fixed-split cross-disease transfer, and
training-pool subsampling (10 seeds) for the adaptation curve.
All feature-wise preprocessing was fit inside training folds only:
standardization, PCA, and any gene selection saw no test cells, and the 70
healthy donors entered no model fit. The primary classifier was logistic
regression ($C=1.0$, lbfgs, balanced class weights); random forest (100 trees,
balanced class weights) is reported as a secondary classifier.
Two safeguards apply to every inferential construction. Any
selection-corrected $P$ must be at least the smallest uncorrected $p$ in its
comparison family, and any nested out-of-fold pooling harness must pass a
unit test, in which a known frozen score refit through the harness must
reproduce its own out-of-fold AUC within 0.02--0.04, before it may produce
significance conclusions (Supplementary Note~\ref{sn:nested_comparison}).

\subsection{Baselines and controls}
\label{ssec:methods_baselines}

Six controls ran under the identical protocol. (i) A from-scratch
PCA + logistic-regression baseline on raw counts: PCA dimensionality was
selected from \{50, 100, 200, 500\} in a pre-registered sweep with
selection performed inside training folds, and we report the 200-component
setting throughout (PCA-200+LR). (ii) A null model using only each patient's
cell count and disease label; within COVID-19, where disease is constant, it
reduces to a cell-count-only model. (iii) A QC-metrics-only model using
per-patient mean UMI counts, detected genes, percent mitochondrial reads, and
cell count, to test whether predictions could be explained by library-quality
or processing artifacts (Supplementary Note~\ref{sn:qc_baseline}).
(iv) A cohort-only model scoring each patient by their cohort's mortality
rate, estimated within training folds only
(Supplementary Note~\ref{sn:cohort_baseline}). (v) A nested incremental-value
family, in which a covariate baseline (QC metrics, cell count, and cohort
identity; L2-regularized logistic regression at $C = 0.3$) is augmented with
each representation score in turn; its pre-specified primary is the single
PCA-200 nested likelihood-ratio test against an empirical permutation null,
chosen because it involves no model selection
(Supplementary Note~\ref{sn:qc_incremental}), and each method is additionally
compared with the null model directly
(Supplementary Note~\ref{sn:method_vs_null}). (vi) An
in-distribution positive control on monocytes from the same source cohorts
(148 patients, 19 fatal), a cell type the models were pretrained on
(Section~\ref{ssec:r4_monocyte}); monocyte identities came from the published
atlas annotations (Ren, Stephenson) or a marker-based call (GSE167363), and a
patient-matched subset analysis is reported in
Supplementary Note~\ref{sn:monocyte_matched}.

\subsection{Cross-disease transfer and adaptation curve}
\label{ssec:methods_transfer}

Cross-disease transfer used fixed disease splits in both directions
(COVID-19$\rightarrow$sepsis and sepsis$\rightarrow$COVID-19) under the primary
severe-as-survived label definition: the forward direction trained on COVID-19
(240 patients) and tested on the 91 GSE279452 sepsis patients, and the reverse
direction trained on all 96 outcome-labeled sepsis patients and tested on
COVID-19, with patient-level scoring and patient-clustered bootstrap
confidence intervals (200 iterations); a
severe-excluded variant is retained as a sensitivity analysis
(Supplementary Note~\ref{sn:severe}). The adaptation curve measured whether
transfer is recovered with target-disease supervision: COVID-trained
classifiers received increasing numbers of sepsis training cells
($n \in \{0, 10, 50, 100, 500, 1000, \text{all}\}$, 10 seeds per size) and were
scored on a held-out 32-patient sepsis test set, a one-third stratified
holdout of the 96-patient labeled sepsis pool.

\subsection{Metrics and calibration}
\label{ssec:methods_metrics}

The primary metric is patient-level out-of-fold ROC-AUC
(Section~\ref{ssec:methods_cv}). Secondary quantities, always named, are the
cell-weighted out-of-fold AUC, the per-fold mean AUC, the leave-one-cohort-out
(LOBO) cohort AUC, and the out-of-fold nested $\Delta$AUC of the
complementarity family (Supplementary Note~\ref{sn:nested_comparison}); no
unqualified ``AUC'' is used elsewhere. Because cohort is
confounded with disease and sampling site in these data, LOBO is treated as a
secondary robustness axis rather than a co-primary estimand; the fixed-split
cross-disease analysis isolates the transfer question directly
(Section~\ref{ssec:methods_transfer}). A LOBO fold is computable only when the
held-out cohort contains at least one fatal and at least one survivor patient;
cohorts failing this rule (34 patients, all in survivor-only cohorts) remain
in the training pool throughout, so pooled LOBO covers 206 of the 240 patients
with all 25 fatal events retained. A fifth named quantity, the
seed-ensemble out-of-fold AUC (ensemble-OOF AUC), averages each patient's
scores across the ten repeats and scores them once; the primary test, all
paired-bootstrap contrasts, fragility refits, score correlations, and the
consensus predictor are computed on these ensemble scores, while the
repeat-mean and repeat-sd report split-to-split stability of the same
protocol.
Calibration was
assessed on patient-level out-of-fold predictions with the Brier score and the
expected calibration error (ECE, 10 equal-width bins). Multiplicity was
controlled within each comparison family with Benjamini--Hochberg FDR, and the
family is stated with every reported $q$.

\subsection{Reproducibility and availability}
\label{ssec:methods_repro}

All analyses ran with fixed seeds (CV seeds 42--51; bootstrap seed 42) and are
fully scripted. Code, per-model environments, the patient-level mapping, and
all result tables are available at [repository URL]. Embeddings were computed
on the UCR HPCC cluster with each model's recommended environment and publicly
released checkpoints (Table~\ref{tab:models}).
